% !TEX root = task-3-4-5.tex
% Editable task bank. Append new BankTask blocks at the end of the document.
\documentclass[12pt,a4paper]{article}
\usepackage{physics-textbook}
\usepackage{graphicx,needspace}
\usetikzlibrary{decorations.pathmorphing}
\newif\ifShowSolutions
\ifdefined\TutorEdition
  \ShowSolutionstrue
\else
  \ShowSolutionsfalse % Change to \ShowSolutionstrue to display solutions.
\fi
\newcounter{banktask}
\newsavebox{\BankTaskBox}
\newcommand{\BankTaskContent}[3]{%
  \noindent\textbf{\thebanktask.}\quad #2\par
  \ifShowSolutions
    \Needspace{50mm}% Keep the solution heading with its opening text or diagram.
    #3\par
    \ifdefined\TutorEdition
      \ifcsname TutorSource#1\endcsname
        \nopagebreak\smallskip{\small\csname TutorSource#1\endcsname\par}%
      \fi
    \fi
  \fi
}
\newcommand{\BankTask}[3]{%
  \refstepcounter{banktask}\par\medskip
  \sbox{\BankTaskBox}{\begin{minipage}{\linewidth}%
    \BankTaskContent{#1}{#2}{#3}\end{minipage}}%
  \ifdim\dimexpr\ht\BankTaskBox+\dp\BankTaskBox\relax>\dimexpr\textheight-20mm\relax
    \Needspace{35mm}\BankTaskContent{#1}{#2}{#3}%
  \else
    \Needspace{\dimexpr\ht\BankTaskBox+\dp\BankTaskBox\relax}%
    \noindent\usebox{\BankTaskBox}\par
  \fi
  \medskip
}
\begin{document}
\BookSubsection{3.4.5}{Потенциал}


% Source sdamgia-1306: https://phys-ege.sdamgia.ru/problem?id=1306
\BankTask{sdamgia-1306}{
Какова разность потенциалов между точками поля, если при перемещении заряда $12 \,\text{мкКл}$ из одной точки в другую электростатическое поле совершает работу $0{,}36 \,\text{мДж}$? (В вольтах.)\par
}{
\smallskip\noindent\textbf{Решение}\par
Разность потенциалов представляет собой работу по переносу единичного заряда между точками поля. Следовательно, разность потенциалов равна $\frac{ 0{,}36\,\text{мДж}}{ 12\,\text{мкКл} } =30\,\text{В}.$\par
Ответ: $30$.\par
}

% Source sdamgia-1429: https://phys-ege.sdamgia.ru/problem?id=1429
\BankTask{sdamgia-1429}{
Модуль напряженности однородного электрического поля равен $100 \,\text{В/м}$. Каков модуль разности потенциалов между двумя точками, расположенными на одной силовой линии поля на расстоянии $5 \,\text{см}$? (В вольтах.)\par
}{
\smallskip\noindent\textbf{Решение}\par
Модуль разности потенциалов между точками, расположенными на одной силовой линии, связана с расстоянием между этими точками и напряженностью однородного электрического поля соотношением\par
\[U=E\Delta x=100\,\text{В/м} \cdot 0{,}05\,\text{м}=5\,\text{В}.\]
Ответ: $5$.\par
}

% Source sdamgia-10644: https://phys-ege.sdamgia.ru/problem?id=10644
\BankTask{sdamgia-10644}{
При перемещении точечного электрического заряда $5 \,\text{мкКл}$ в электростатическом поле из точки $1$ в точку $2$ действующая со стороны этого поля сила совершает работу $17 \,\text{мкДж}$. При перемещении того же заряда из точки $1$ в точку $3$ в этом же электростатическом поле действующая со стороны поля сила совершает работу $7 \,\text{мкДж}$. Чему равна разность потенциалов между точками $3$ и $2$ этого поля?\par
}{
\smallskip\noindent\textbf{Решение}\par
Потенциальная энергия заряда $q$ в электростатическом поле в некоторой точке связана с величиной потенциала $\varphi$ в этой точке соотношением: $W=q\varphi.$ Таким образом, изменение потенциальной энергии заряда при перемещении из точки $1$ в точку $2$ равно:\par
\[W_{12}=W_{1} - W_{2}=q ( \varphi_{1} - \varphi_{2} ) .\]
При перемещении заряда из точки $1$ в точку $3$ совершается работа\par
\[W_{13}=q ( \varphi_{1} - \varphi_{3} ) .\]
Разность потенциалов между точками $2$ и $3$ составляет\par
\[\varphi_{3} - \varphi_{2}= \frac{ W_{12} - W_{13}}{ q } =2\,\text{В}.\]
Ответ: $2 \,\text{В}$.\par
}

% Source sdamgia-4104: https://phys-ege.sdamgia.ru/problem?id=4104
\BankTask{sdamgia-4104}{
Двум металлическим пластинам площадью $S$ каждая сообщили равные по модулю, но противоположные по знаку заряды +$Q$ и -$Q$. Пластины расположили на малом расстоянии $d$ друг от друга. Установите соответствие между физическими величинами и формулами, по которым их можно рассчитать.\par

\begin{center}\begin{minipage}[t]{.47\linewidth}\textbf{Формулы}\par\smallskip $A$) $\frac{ Qd}{ \varepsilon_{0} S }$\par
Б) $\frac{ Q ^2 d}{ 2\varepsilon_{0} S }$\end{minipage}\hfill\begin{minipage}[t]{.47\linewidth}\textbf{Физические величины}\par\smallskip $1$. Напряженность электрического поля между пластинами\par
$2$. Разность потенциалов между пластинами\par
$3$. Емкость системы, состоящей из двух таких пластин\par
$4$. Энергия электрического поля, заключенного между этими пластинами\end{minipage}\end{center}
\par
\begin{center}\small\setlength{\tabcolsep}{3pt}\begin{tabular}{|p{0.45\linewidth}|p{0.45\linewidth}|}\hline
$A$ & Б\\\hline
 & \\\hline
\end{tabular}\end{center}\par
}{
\smallskip\noindent\textbf{Решение}\par
Описанная в задаче система представляет собой заряженный плоский конденсатор. Емкость такого конденсатора равна $C= \frac{ \varepsilon_{0} S}{ d } .$ Энергия плоского конденсатора дается выражением: $W= \frac{ Q ^2 }{ 2C } = \frac{ Q ^2 d}{ 2\varepsilon_{0} S } ,$ следовательно, формула Б описывает энергию электрического поля, заключенного между пластинами. Разность потенциалов между пластинами равна $U= \frac{ Q}{ C } = \frac{ Qd}{ \varepsilon_{0} S } .$ Следовательно, формуле А соответствует пункт $2$. Наконец, чтобы завершить разговор по поводу данной задачи, напомним выражение для напряженности электрического поля между пластинами: $E= \frac{ U}{ d } = \frac{ Q}{ \varepsilon_{0} S } .$\par
Ответ: $24$.\par
}

% Source sdamgia-9029: https://phys-ege.sdamgia.ru/problem?id=9029
\BankTask{sdamgia-9029}{
Однородное электростатическое поле создано равномерно заряженной протяженной горизонтальной пластиной. Линии напряженности поля направлены вертикально вверх (см. рис.).\par
$1$. Если в точку А поместить пробный точечный отрицательный заряд, то на него со стороны пластины будет действовать сила, направленная вертикально вниз.\par
$2$. Пластина имеет отрицательный заряд.\par
$3$. Потенциал электростатического поля в точке В ниже, чем в точке С.\par
$4$. Напряженность поля в точке А меньше, чем в точке С.\par
$5$. Работа электростатического поля по перемещению пробного точечного отрицательного заряда из точки А и в точку В равна нулю.\par

\begin{center}\begin{tikzpicture}[book diagram,>=stealth]
\draw[thick](-.3,0)--(3.4,0);\foreach\x in{0,1,2,3}\draw[->](\x,0)--(\x,3.1);\node[right]at(2.2,2.65){$\vec E$};\draw(0,1)--(3,1)--(3,2.5)--cycle;\node[left]at(0,1){$A$};\node[right]at(3,1){$B$};\node[right]at(3,2.5){$C$};\draw(2.8,1)--(2.8,1.2)--(3,1.2);
\end{tikzpicture}\end{center}
\par
}{
\smallskip\noindent\textbf{Решение}\par
Поскольку линии напряженности выходят из пластины, то, значит, пластина заряжена положительно (утверждение $2$ неверно). Если в точку А поместить пробный точечный отрицательный заряд, то на него будет действовать сила, направленная вертикально вниз (утверждение $1$ верно).\par
Потенциал убывает при удалении от положительной пластины (утверждение $3$ неверно).\par
Напряженность поля одинакова во всех точках над пластиной (утверждение $4$ неверно).\par
Точки А и $B$ — это точки одинакового потенциала, а значит, работа по перемещению пробного заряда между этими точками равна нулю (утверждение $5$ верно).\par
Ответ: $15$.\par
}

% Source sdamgia-10474: https://phys-ege.sdamgia.ru/problem?id=10474
\BankTask{sdamgia-10474}{
Две параллельные металлические пластины больших размеров расположены на расстоянии $d$ друг от друга и подключены к источнику постоянного напряжения (рис. $1$). Пластины закрепили на изолирующих подставках и спустя длительное время отключили от источника (рис. $2$).\par
$1$. Напряженность электрического поля в точке А больше, чем в точке В.\par
$2$. Потенциал электрического поля в точке А больше, чем в точке С.\par
$3$. Если увеличить расстояние между пластинами $d$, то напряженность электрического поля в точке С не изменится.\par
$4$. Если уменьшить расстояние между пластинами $d$, то заряд правой пластины не изменится.\par
$5$. Если пластины полностью погрузить в керосин, то энергия электрического поля конденсатора останется неизменной.\par

\begin{center}\begin{tikzpicture}[book diagram,>=stealth]
\foreach\shift in{0,4}{\begin{scope}[xshift=\shift cm]\draw[thick](0,0)--(0,2.6)(2,0)--(2,2.6);\draw[<->](0,3)--(2,3)node[midway,above]{$d$};\draw(0,2.6)--(0,3.2)(2,2.6)--(2,3.2);\end{scope}}\draw(0,1.3)--(-.8,1.3)--(-.8,-.8)--(.65,-.8);\draw(.65,-.8)circle(.1);\draw(1.35,-.8)circle(.1);\draw(1.35,-.8)--(2.8,-.8)--(2.8,1.3)--(2,1.3);\node[above]at(.65,-.7){$-$};\node[above]at(1.35,-.7){$+$};\node at(1,-1.3){Рис.\ 1};\draw(4,1.3)--(3.5,1.3)(6,1.3)--(6.5,1.3);\foreach\x/\y/\s in{4.55/.55/A,4.55/1.7/B,5.3/1.15/C}{\fill(\x,\y)circle(1.5pt);\node[above]at(\x,\y){$\s$};}\node at(5,-1.3){Рис.\ 2};
\end{tikzpicture}\end{center}
\par
}{
\smallskip\noindent\textbf{Решение}\par
После того как длительное время пластины были подключены к источнику постоянного напряжения, они зарядились: левая пластина — отрицательно, правая — положительно.\par
$1$. Внутри плоского заряженного конденсатора электрическое поле однородно. Напряженности поля в точках А и В одинаковые. Утверждение $1$ неверно.\par
$2$. Потенциал электрического поля внутри конденсатора убывает от положительной пластины к отрицательной. Потенциал электрического поля в точке А меньше, чем в точке С. Утверждение $2$ неверно.\par
$3$. Поскольку пластины отключены от источника, то заряд и его поверхностная плотность на них не меняется при изменении расстояния. Значит, не будет изменяться и напряженность электрического поля между пластинами. Утверждение $3$ верно.\par
$4$. Заряд пластин остается постоянным, независимо от того, сдвигают пластины или нет. Утверждение $4$ верно.\par
$5$. Диэлектрическая проницаемость керосина больше $1$. При полном погружении в керосин энергия электрического поля конденсатора $W= \frac{ q ^2 }{ 2C }$ уменьшится. Утверждение $5$ неверно.\par
Ответ: $34$.\par
}

% Source sdamgia-6359: https://phys-ege.sdamgia.ru/problem?id=6359
\BankTask{sdamgia-6359}{
Сплошной металлический шар радиусом $R$, находящийся в вакууме, имеет заряд $Q$. Установите соответствие между физическими величинами и формулами, по которым их можно рассчитать.\par

\begin{center}\begin{minipage}[t]{.47\linewidth}\textbf{Физические величины}\par\smallskip А) модуль напряженности электрического поля на расстоянии $2R$ от центра шара\par
Б) потенциал поверхности шара\end{minipage}\hfill\begin{minipage}[t]{.47\linewidth}\textbf{Формулы}\par\smallskip $1$) $\frac{ 1}{ 4 \pi \varepsilon_{0} } \cdot \frac{ Q}{ 2R }$ \par
$2$) $0$\par
$3$) $\frac{ 1}{ 4 \pi \varepsilon_{0} } \cdot \frac{ Q}{ R }$ \par
$4$) $\frac{ 1}{ 4 \pi \varepsilon_{0} } \cdot \frac{ Q}{ 4R ^2 }$\end{minipage}\end{center}
\par
\begin{center}\small\setlength{\tabcolsep}{3pt}\begin{tabular}{|p{0.45\linewidth}|p{0.45\linewidth}|}\hline
$A$ & Б\\\hline
 & \\\hline
\end{tabular}\end{center}\par
}{
\smallskip\noindent\textbf{Решение}\par
А) Модуль напряженности сферически симметричного распределения заряда вне этого заряда на расстоянии $r$ рассчитывается по формуле $E=k\frac{Q}{r^2}=\frac1{4\pi\varepsilon_0}\cdot\frac Q{r^2}.$ Следовательно, на расстоянии $2R$ модуль напряженности $E= \frac{ 1}{ 4 \pi \varepsilon _{0} } \cdot \frac{ Q}{ 4R ^2 } .$\par
Б) Потенциал по сферически симметричного распределения заряда вне этого заряда на расстоянии $r$ рассчитывается по формуле $\varphi=k \frac{ Q}{ r } = \frac{ 1}{ 4 \pi \varepsilon _{0} } \cdot \frac{ Q}{ r } .$ Значит, на поверхности шара, то есть на расстоянии $R$, потенциал $\varphi= \frac{ 1}{ 4 \pi \varepsilon _{0} } \cdot \frac{ Q}{ R } .$\par
Ответ: $43$.\par
}

% Source sdamgia-8876: https://phys-ege.sdamgia.ru/problem?id=8876
\BankTask{sdamgia-8876}{
Частица массой $0{,}08 \,\text{мг}$, имеющая заряд $10^{-10}\,\text{Кл}$, покоится в точке А. При включении горизонтального однородного электрического поля эта частица, двигаясь по горизонтали вдоль силовой линии, смещается в точку В. Напряжение между точками $A$ и $B$ равно $1 \,\text{В}$. Чему равна скорость частицы в точке В? (В метрах в секунду.)\par
}{
\smallskip\noindent\textbf{Решение}\par
Электрическое поле совершает работу по перемещению заряда: $A=qU.$\par
Работа по перемещению заряда равна изменению кинетической энергии $\Delta W_{\text{к}}= \frac{ m v ^2 }{ 2 } - \frac{ m v _{0} ^2 }{ 2 } ,$ причем, покоящаяся вначале частица кинетической энергией не обладала. Таким образом, получаем:\par
\[qU= \frac{ m v ^2 }{ 2 } .\]
Отсюда скорость частицы равна:\par
\[v=\sqrt{\frac{2qU}m}=\sqrt{\frac{2\cdot10^{-10}\cdot1}{0{,}08\cdot10^{-6}}}=0{,}05\,\text{м/с}.\]
Ответ: $0{,}05 \,\text{м/с}$.\par
}

% Source sdamgia-2990: https://phys-ege.sdamgia.ru/problem?id=2990
\BankTask{sdamgia-2990}{
Пылинка, имеющая массу $10 ^{ - 8 } \,\text{г}$ и заряд $( - 1{,}8 ) \cdot 10 ^{ - 14 } \,\text{Кл},$ влетает в электрическое поле вертикального плоского конденсатора в точке, находящейся посередине между его пластинами (см. рис., вид сверху).\par
Чему должна быть равна минимальная скорость, с которой пылинка влетает в конденсатор, чтобы она смогла пролететь его насквозь? Длина пластин конденсатора $10 \,\text{см}$, расстояние между пластинами $1 \,\text{см}$, напряжение на пластинах конденсатора $5000 \,\text{В}$. Система находится в вакууме. Действием силы тяжести пренебречь.\par

\begin{center}\begin{tikzpicture}[book diagram,>=stealth]
\draw(0,1.1)rectangle(5,1.4);\draw(0,-1.1)rectangle(5,-1.4);\foreach\x in{.4,1.4,...,4.4}{\node at(\x,1.25){$+$};\node at(\x,-1.25){$-$};}\draw[thin](.2,0)--(5,0);\draw[fill=white](.15,0)circle(.12);\draw[->,thick](.27,0)--(1.2,0)node[midway,above]{$\vec v$};
\end{tikzpicture}\end{center}
\par
}{
\smallskip\noindent\textbf{Решение}\par
Сила, действующая на частицу в конденсаторе со стороны поля: $F_{\text{эл}}=E|q|.$ Связь напряженности электрического поля с напряжением на пластинах конденсатора: $E= \frac{ U}{ d } ,$ где $d$ — расстояние между пластинами. Второй закон Ньютона в проекции на ось, перпендикулярную пластинам: $F_{\text{эл}}=ma,$ или $\frac{ |q|U}{ d } =ma.$\par
Рассмотрим движение частицы в электрическом поле конденсатора, выбрав оси: Ox вдоль пластины, Oy перпендикулярно вверх на рисунке. Так как пылинка отрицательная, то вектор электрической силы противоположно направлен с вектором напряженности (вверх), по второму закону Ньютона вектор ускорения сонаправлен с вектором силы. Таким образом, движение в поле конденсатора по оси Ox равномерное со скоростью $\upsilon$, перемещение $s_{x}$ равно длине пластин $l$. По оси Oy движение равноускоренное, причем проекция начальной скорости на ось Oy равна $0$, перемещение — $s_{y}.$\par
Сила со стороны электрического поля действует в вертикальном направлении, в вертикальной плоскости будет движение под действием этой силы, по параболе. Время полета вдоль пластин составляет $t= \frac{ l}{ v } .$ За это время частица сместится в сторону пластины на $s_{y}= \frac{ at ^2 }{ 2 } = \frac{ |q|Ul ^2 }{ 2dm v ^2 } .$ Условием пролета насквозь является $s \le \frac{ d}{ 2 } .$ Откуда получаем минимальную скорость:\par
\[v=\frac ld\sqrt{\frac{|q|U}{m}}=\frac{0{,}1}{0{,}01}\sqrt{\frac{1{,}8\cdot10^{-14}\cdot5000}{10^{-11}}}=30\,\text{м/с}.\]
Ответ: $v =30\,\text{м/с}.$\par
}

% Source sdamgia-2993: https://phys-ege.sdamgia.ru/problem?id=2993
\BankTask{sdamgia-2993}{
Маленький шарик с зарядом $q=4 \cdot 10 ^{ - 8 } \,\text{Кл}$ и массой $3 \,\text{г}$, подвешенный на невесомой нити с коэффициентом упругости $100 \,\text{Н/м}$, находится между вертикальными пластинами плоского воздушного конденсатора. Расстояние между обкладками конденсатора $5 \,\text{см}$. Какова разность потенциалов между обкладками конденсатора, если удлинение нити $0{,}5 \,\text{мм}$?\par

\begin{center}\begin{tikzpicture}[book diagram,>=stealth]
\draw[thick](-1.8,-1)--(-1.8,2.3)(1.1,-1)--(1.1,2.3);\draw(-1.4,2.4)--(.8,2.4);\foreach\x in{-1.4,-1.1,...,.7}\draw(\x,2.4)--++(.2,.2);\draw(0,2.4)--(-1.3,.2);\draw[fill=white](-1.3,.2)circle(.2);
\end{tikzpicture}\end{center}
\par
}{
\smallskip\noindent\textbf{Решение}\par
Изобразим силы, действующие на шарик:\par

\begin{center}\begin{tikzpicture}[book diagram,>=stealth]
\draw[thick](-1.8,-1)--(-1.8,2.3)(1.1,-1)--(1.1,2.3);\draw(-1.4,2.4)--(.8,2.4);\foreach\x in{-1.4,-1.1,...,.7}\draw(\x,2.4)--++(.2,.2);\draw(0,2.4)--(-1.3,.2);\draw[fill=white](-1.3,.2)circle(.2) (0,0)circle(.2);\node[above left]at(-1.3,.2){$+q$};\draw[->](0,2.4)--(-2.3,2.4)node[below]{$x$};\draw[->](0,2.4)--(0,-1.25)node[below]{$y$};\draw[->,thick](-1.3,.2)--(-1.3,-.7)node[right]{$m\vec g$};\draw[->,thick](-1.3,.2)--(-2.3,.2)node[above]{$\vec F_{\text{э}}$};\draw[->,thick](-1.3,.2)--(-.65,1.3)node[left]{$\vec F_{\text{упр}}$};\draw(0,1.65)arc(-90:-120:.75);\node at(-.22,1.48){$\alpha$};
\end{tikzpicture}\end{center}
\par
Шарик находится в равновесии, если выполняется условие $\vec{F}_{\text{упр}} + m \vec{g} + \vec{F}_{\text{э}}=0.$ По закону Гука модуль силы упругости $F_{\text{упр}}=k\Delta l.$ Сила, с которой электрическое поле действует на заряженное тело, равна по модулю $F_{\text{э}}=Eq$ и сонаправлена с вектором напряженности, так как заряд тела положительный.\par
Проекция условия равновесия на координатные оси:\par
Ox: $F_{\text{упр}} \cdot \sin \alpha =F_{\text{э}},$\par
Oy: $F_{\text{упр}} \cdot \cos \alpha =mg.$\par
Объединяя уравнения, получаем:\par
\[k\Delta l \cdot \sin \alpha =qE,\]
\[k\Delta l \cdot \cos \alpha =mg.\]
Возведем оба равенства в квадрат и сложим их: $( k\Delta l ) ^2 = ( mg ) ^2 + ( qE ) ^2 ,$ откуда:\par
\[E= \frac{ \sqrt{( k\Delta l ) ^2 - ( mg ) ^2} }{ q } .\]
Напряженность электрического поля в конденсаторе: $E= \frac{ U}{ d } .$\par
Таким образом, напряжение между обкладками конденсатора:\par
\[\begin{aligned}U&=\frac{d\sqrt{(k\Delta l)^2-(mg)^2}}q\\&=\frac{5\cdot10^{-2}\sqrt{(100\cdot5\cdot10^{-4})^2-(3\cdot10^{-3}\cdot10)^2}}{4\cdot10^{-8}}=50\,\text{кВ}.\end{aligned}\]
Ответ: $U=50\,\text{кВ}.$\par
}

\end{document}
